\PassOptionsToPackage{dvispnames, table}{xcolor}
\documentclass[aspectratio=169,hyperref={colorlinks=true}]{beamer}
\usepackage[beamer]{prettytex/base}
\usepackage{prettytex/math}
\usepackage{prettytex/gfx}
\usepackage[backend=biber,citestyle=authoryear,bibstyle=numeric,hyperref=true,maxbibnames=99]{biblatex}
\usepackage{csquotes}
\usepackage{fontawesome}
\addbibresource{../literature/sources.bib}
\usetheme{oxonian}

% \definecolor{applegreen}{rgb}{0.55,0.71,0.0}
% \definecolor{americanrose}{rgb}{1.0,0.01,0.24}
% \definecolor{aqua}{rgb}{0.0,1.0,1.0}
% TODO: add name to footline

\title{General Kernel Spectral Methods for Equilibrium Measures \\ \normalsize MMSC Dissertation}
\titlegraphic{\includegraphics[width=2cm]{theme/logos/ox-circle.eps}}
\author{Peter Julius Waldert}
\institute{Mathematical Institute \\ University of Oxford}
\date{} %\today

\begin{document}
  {\setbeamertemplate{footline}{}\frame{\titlepage}}

  \section{Intro}
  \begin{frame}{Project Motivation and Goal}
    \begin{itemize}
      \item Simulate Many-Particle-Systems and find their Equilibrium Distribution.
      \item Interactions through (power-law) Attraction-Repulsion Potentials
            $$K(r) = \frac{r^\alpha}{\alpha} - \frac{r^\beta}{\beta}\,,\quad \alpha, \beta \in \R\,.$$
      \item Each particle $i=1, ..., N$ at position $\vec{x_i}$ and time $t$ follows
            $$\frac{\dd^2 \vec{x_i}}{\ddt^2} = f\left(\norm{\frac{\dd \vec{x_i}}{\ddt}}_2\right) \frac{\dd\vec{x_i}}{\ddt} - \frac{1}{N} \sum_{j=1, i\neq j}^{N} \nabla K\left(\norm{\vec{x_i} - \vec{x_j}}_2\right)\,,$$
            \parencite{2020-power-law-kernels, 2021-arbitrary-dimensions}.
      \item Add tikzpicture of particle box with interaction
      \item If repulsive part is larger (so $\beta > \alpha$), there is no equilibium distribution as particles simply continue repelling one another out to infinity
    \end{itemize}
  \end{frame}

  \begin{frame}{An Integral Equation}
    The total energy is $$E = \sum_{i=1}^{N} \sum_{j=1, j \neq i}^{N} K\left(\norm{\vec{x_i} - \vec{x_j}}_2\right)\,,$$
    which can be equivalently expressed as
    $$E = \frac{1}{2} \iint K\left(\norm{\vec{x} - \vec{y}}_2\right) \,\dd\rho(\vec{x})\,\dd\rho(\vec{y})\,,$$ where $\dd\rho = \dd\rho(x)\ddx$ is a measure chosen such that $$M = \int \dd\rho = \int_{\supp(\rho)} \rho(x) \,\ddx = 1\,.$$
  \end{frame}

  \begin{frame}{Simulation Results}
    What we expect:

    [Plots of simulator equilibium measure]
  \end{frame}

  \begin{frame}{Spectral Methods}
    \begin{itemize}
      \item In order to find $\rho(x)$, we consider the following Ansatz
            $$\rho(\vec{x}) = \sum_{n=0}^{\infty} f_n p_n(\vec{x})$$
            by which we construct a spectral method for the solution of the above integral equation.
      \item Using the Jacobi polynomial basis.
      \item Minimization routine of $E$ over coefficients in $\rho$, as a subroutine of outer minimisation over $a$ and $b$ the bounds of the box (simpler: $a = b = r$).
      \item We see that we do not need the inner optimisation routine, because TODO.
    \end{itemize}
  \end{frame}

  \begin{frame}{Results from the Spectral Solver}
    Here are some plots

    If too ambitious, just use those from paper
  \end{frame}

  \begin{frame}{Project Goals}
    \begin{itemize}
      \item Establish theory behind spectral methods in the Jacobi basis.
      \item Use kernel expansions to construct an equilibrium measure method on $[-1, 1]$.
      \item Consider the $d$-dimensional unit ball.
      \item Numerically solve for $\rho(x)$ using a fully implemented spectral solver.
      \item Implement a particle simulator for the same kernel $K$.
      \item Compare results of the two methods.
      \item First numerical method on continuous space problems!! (we are the fastest anyway!).
      \item Optional: Potentially, extend to l-Morse potentials.
      \item Spread Open Source!
    \end{itemize}
  \end{frame}

  \begin{frame}{}
    Questions?
  \end{frame}

  \begin{frame}[allowframebreaks]
    \frametitle{Bibliography}
    \nocite{*}
    \printbibliography[heading=bibnumbered]
  \end{frame}
\end{document}
